\documentclass[11pt,a4paper]{article}
\usepackage[utf8]{inputenc}
\usepackage[T1]{fontenc}
\usepackage[portuguese]{babel}
\usepackage{geometry}
\geometry{margin=2.5cm}
\usepackage{listings}
\usepackage{xcolor}
\usepackage{amsmath}
\usepackage{amssymb}
\usepackage{hyperref}
\usepackage{enumitem}
\usepackage{tcolorbox}
\usepackage{tabularx}
\usepackage{booktabs}

\definecolor{codebg}{rgb}{0.95,0.95,0.95}
\definecolor{codegreen}{rgb}{0.1,0.5,0.1}
\definecolor{codepurple}{rgb}{0.5,0.1,0.5}
\definecolor{codegray}{rgb}{0.5,0.5,0.5}

\lstdefinestyle{python}{
    backgroundcolor=\color{codebg},
    commentstyle=\color{codegray}\itshape,
    keywordstyle=\color{codepurple}\bfseries,
    stringstyle=\color{codegreen},
    basicstyle=\ttfamily\small,
    breaklines=true,
    frame=single,
    language=Python,
    showstringspaces=false,
    tabsize=4,
    literate={á}{{\'a}}1 {ã}{{\~a}}1 {é}{{\'e}}1 {í}{{\'i}}1 {ó}{{\'o}}1 {õ}{{\~o}}1 {ú}{{\'u}}1 {ç}{{\c{c}}}1 {Á}{{\'A}}1 {Ã}{{\~A}}1 {É}{{\'E}}1 {Í}{{\'I}}1 {Ó}{{\'O}}1 {Õ}{{\~O}}1 {Ú}{{\'U}}1 {Ç}{{\c{C}}}1 {â}{{\^a}}1 {ê}{{\^e}}1 {ô}{{\^o}}1 {û}{{\^u}}1 {Â}{{\^A}}1 {Ê}{{\^E}}1 {Ô}{{\^O}}1 {Û}{{\^U}}1 {à}{{\`a}}1 {è}{{\`e}}1 {ì}{{\`i}}1 {ò}{{\`o}}1 {ù}{{\`u}}1
}
\lstset{style=python}

\newtcolorbox{solucao}{
  colback=yellow!5!white,
  colframe=orange!60!black,
  title=Solução proposta,
  fonttitle=\bfseries
}

\title{\textbf{Data Analysis with Python 2024--2025}\\Parte 4: Pandas (Parte 1)\\\large Soluções dos Exercícios}
\author{Soluções independentes baseadas nos enunciados originais}
\date{}

\begin{document}
\maketitle

\section*{Nota Importante}
Este material é uma adaptação das soluções independentes para os exercícios baseados no curso \textit{Data Analysis with Python 2024-2025}, oferecido pela \textbf{The University of Helsinki MOOC Center}. As soluções abaixo não são o gabarito oficial do curso.

\tableofcontents
\newpage

\section{Soluções - Capítulo 3 (Pandas Parte 1)}

\subsection*{Exercício 3.13 -- Ler série (\textit{read series})}
\begin{solucao}
\begin{lstlisting}
import pandas as pd

def read_series():
    indices = []
    values = []
    while True:
        line = input("Digite o indice e o valor (ou deixe em branco para sair): ")
        if not line:
            break
        parts = line.split()
        if len(parts) != 2:
            raise Exception("Malformed input")
        indices.append(parts[0])
        values.append(parts[1])
    return pd.Series(values, index=indices)

def main():
    try:
        s = read_series()
        print(s)
    except Exception as e:
        print(f"Erro: {e}")

if __name__ == "__main__":
    main()
\end{lstlisting}
\end{solucao}

\subsection*{Exercício 3.14 -- Operações em séries (\textit{operations on series})}
\begin{solucao}
\begin{lstlisting}
import pandas as pd

def create_series(L1, L2):
    s1 = pd.Series(L1, index=['a', 'b', 'c'])
    s2 = pd.Series(L2, index=['a', 'b', 'c'])
    return s1, s2

def modify_series(s1, s2):
    s1['d'] = s2['b']
    del s2['b']
    return s1, s2

def main():
    s1, s2 = create_series([1, 2, 3], [4, 5, 6])
    print("S1 e S2 Originais:")
    print(s1)
    print(s2)
    
    s1, s2 = modify_series(s1, s2)
    print("\nApós modificação:")
    print(s1)
    print(s2)
    
    print("\nSoma de s1 e s2 (alinhado pelos índices):")
    print(s1 + s2)

if __name__ == "__main__":
    main()
\end{lstlisting}
\end{solucao}

\subsection*{Exercício 3.15 -- Série inversa (\textit{inverse series})}
\begin{solucao}
\begin{lstlisting}
import pandas as pd

def inverse_series(s):
    return pd.Series(s.index, index=s.values)

def main():
    s = pd.Series([1, 2, 2], index=['a', 'b', 'c'])
    print("Serie original:")
    print(s)
    
    inv = inverse_series(s)
    print("\nSerie invertida:")
    print(inv)
    
    print("\nAcessando o indice repetido '2':")
    # Retornara todos os elementos associados a este indice
    print(inv[2]) 

if __name__ == "__main__":
    main()
\end{lstlisting}

\textbf{Resposta para a pergunta do exercício:} Quando um valor aparece várias vezes na Série original, esse valor se tornará um índice repetido na nova Série (Série inversa). Se você usar esse valor repetido para indexar a nova Série (\texttt{inv[valor]}), o Pandas retornará uma nova Série contendo todos os elementos que possuem aquele índice (em vez de retornar apenas um único valor escalar).
\end{solucao}

\vspace{2em}
\noindent\textbf{Créditos:} Este conteúdo foi originalmente criado pela \textit{The University of Helsinki MOOC Center} para o curso \textit{Data Analysis with Python}.
\end{document}
